Neuro-symbolic (NeSy) systems promise to combine neural perception with symbolic reasoning,
yielding data efficiency, interpretable symbols, and compositional reuse. But these promises rest
on a hidden assumption: that the system actually \emph{grounds} its latent symbols correctly. We
ask, on the canonical \mnistadd benchmark, when wiring a fixed symbolic ``$+$'' into a shared neural
digit classifier---trained only on the \emph{sum} of two images, never on the digits themselves---
produces genuine perception-plus-reasoning, and when it only appears to. Across 210 training runs
(5 seeds per configuration), we report four findings measured side by side with confidence intervals
and significance tests. First, integration is a large data-efficiency win: at 1{,}000 training pairs
the NeSy model reaches $85.4\%$ sum accuracy versus $22.0\%$ for a pure-neural baseline (Welch
$p{=}2.9{\times}10^{-13}$), matching a fully digit-supervised oracle. Second, high task accuracy can
hide a broken model: changing only the symbolic operator to $(a{+}b)\bmod 10$ makes correct grounding
\emph{never} emerge---raw grounding $\approx0.1\%$ among seeds that learn the task---because the
network reads digits via a systematic relabeling that leaves the modular sum invariant, a reasoning
shortcut visible only when grounding is measured permutation-aligned. Third, just $10$ labeled images
(one per class) eliminate the shortcut and stabilize training (grounding $\approx0.1\%\to97.7\%$),
whereas unsupervised regularizers do not. Fourth, correct grounding unlocks zero-shot compositional
generalization to $2$- and $3$-digit addition ($95.7\%$/$93.7\%$) that the neural baseline structurally
cannot achieve ($0\%$). The lesson for practitioners: task accuracy is not enough; grounding must be
measured, the operator's symmetry determines whether symbols are identifiable, and a pinch of symbol
supervision is a cheap, decisive fix.
